```latex
\documentclass{article}
\usepackage{amsmath,amsthm,amssymb}

\newtheorem{lemma}{Lemma}
\newtheorem{proposition}{Proposition}
\newtheorem{remark}{Remark}

\title{On tame almost complex structures for symplectic simple normal crossings divisors}
\author{}
\date{}

\begin{document}
\maketitle

\begin{remark}
In what follows, a \emph{symplectic simple normal crossings divisor} is understood in the weak sense: each irreducible component is a symplectic submanifold, and all intersections are transverse and symplectic. If one instead includes $\omega$-orthogonality of the branches in the definition, then the answer to the question is immediately yes via the standard Weinstein normal form and the standard complex structure.
\end{remark}

\noindent
The answer is \textbf{no}. We give a counterexample on $\mathbb{T}^4$.

\section{The counterexample}

Let $M=\mathbb{T}^4=\mathbb{R}^4/\mathbb{Z}^4$ with coordinates
$x_1,y_1,x_2,y_2$. Fix $c>1$ and define
\[
\omega_c=
dx_1\wedge dy_1+dx_2\wedge dy_2
+c\,dx_1\wedge dx_2+c\,dy_1\wedge dy_2.
\]
Since
\[
\omega_c^2
=
2(c^2-1)\,dx_1\wedge dx_2\wedge dy_1\wedge dy_2
\neq 0,
\]
$\omega_c$ is symplectic.

Define the codimension-two symplectic tori
\[
D_1=\{x_2=y_2=0\},\qquad D_2=\{x_1=y_1=0\},
\]
and put $D=D_1\cup D_2$. The restrictions are
\[
\omega_c|_{D_1}=dx_1\wedge dy_1,\qquad
\omega_c|_{D_2}=dx_2\wedge dy_2,
\]
so $D_1$ and $D_2$ are symplectic. They intersect transversely at the single point
\[
p=(0,0,0,0),
\]
and a point is symplectic. Thus $D$ is a symplectic simple normal crossings divisor.

At $p$, let
\[
W_1=T_pD_1=\operatorname{span}(\partial_{x_1},\partial_{y_1}),
\qquad
W_2=T_pD_2=\operatorname{span}(\partial_{x_2},\partial_{y_2}).
\]
In the ordered basis
$(e_1,f_1,e_2,f_2)=(\partial_{x_1},\partial_{y_1},\partial_{x_2},\partial_{y_2})$,
the matrix of $\omega_c$ is
\[
\begin{pmatrix}
0&1&c&0\\
-1&0&0&c\\
-c&0&0&1\\
0&-c&-1&0
\end{pmatrix}.
\]
Its Pfaffian is
\[
\operatorname{Pf}(\omega_c)=1-c^2<0.
\]

\section{A linear obstruction}

\begin{lemma}
Let $(V,\omega)$ be a $4$-dimensional symplectic vector space, and let
$V=W_1\oplus W_2$ be a direct sum of symplectic $2$-planes.
Suppose that for some bases $e_1,f_1\in W_1$ and $e_2,f_2\in W_2$ with
$\omega(e_i,f_i)=1$, the Pfaffian of $\omega$ in the basis
$(e_1,f_1,e_2,f_2)$ is negative.
Then there is no $\omega$-tame almost complex structure $J$ on $V$
such that $J(W_i)\subseteq W_i$ for $i=1,2$.
\end{lemma}

\begin{proof}
Write the cross matrix $A$ by
\[
A=\begin{pmatrix}
\omega(e_1,e_2)&\omega(e_1,f_2)\\
\omega(f_1,e_2)&\omega(f_1,f_2)
\end{pmatrix}.
\]
The Pfaffian of $\omega$ in the given basis is
\[
\operatorname{Pf}(\omega)=1-\det A.
\]
Thus $\det A>1$. Since $\det A>0$, we can write
$A=cS$ with $c=\sqrt{\det A}>1$ and $S\in SL(2,\mathbb{R})$.
Replacing the basis of $W_2$ by its image under $S^{-1}$, which is a symplectic change of basis of $W_2$, transforms $A$ into $cI$.
Hence we may assume
\[
\omega(e_1,e_2)=c,\qquad \omega(f_1,f_2)=c,
\]
and all other cross pairings between $W_1$ and $W_2$ are zero.

Assume that an $\omega$-tame $J$ with $J(W_i)\subseteq W_i$ exists.
Because $J$ preserves $W_1$, we may write
\[
Je_1=\alpha e_1+a f_1
\]
with
\[
a=\omega(e_1,Je_1)>0.
\]
Similarly, since $J$ preserves $W_2$, write
\[
Jf_2=\gamma e_2+\delta f_2.
\]
Tameness gives
\[
b:=\omega(f_2,Jf_2)=-\gamma>0,
\]
so $\gamma=-b$.

Let $q$ be the positive definite quadratic form
\[
q(v)=\omega(v,Jv),
\]
and let its symmetric bilinear part also be denoted by $q$.
We compute
\[
\omega(e_1,Jf_2)
=
\omega(e_1,\gamma e_2+\delta f_2)
=
\gamma c
=
-bc,
\]
and
\[
\omega(f_2,Je_1)
=
\omega(f_2,\alpha e_1+a f_1)
=
a\,\omega(f_2,f_1)
=
-ac.
\]
Therefore
\[
q(e_1,f_2)
=
\frac12\bigl(\omega(e_1,Jf_2)+\omega(f_2,Je_1)\bigr)
=
-\frac c2(a+b).
\]
Also
\[
q(e_1,e_1)=a,\qquad q(f_2,f_2)=b.
\]
Thus the restriction of $q$ to the $2$-plane spanned by $e_1$ and $f_2$ has matrix
\[
\begin{pmatrix}
a&-\frac c2(a+b)\\[2mm]
-\frac c2(a+b)&b
\end{pmatrix}.
\]
Its determinant is
\[
ab-\frac{c^2}{4}(a+b)^2.
\]
But $a,b>0$, so $(a+b)^2\ge 4ab$, and since $c>1$,
\[
ab-\frac{c^2}{4}(a+b)^2
\le
ab-c^2ab
=
(1-c^2)ab
<0.
\]
This contradicts the positive definiteness of $q$. Hence no such $J$ exists.
\end{proof}

\section{The obstruction persists under the allowed deformations}

The question allows replacing $\omega$ by a symplectic form connected to it
through a path of symplectic forms, all cohomologous to $\omega$, along which
$D$ remains a symplectic simple normal crossings divisor. For the example above,
the linear obstruction at $p$ cannot be removed by such a deformation.

Indeed, let $\omega_t$ be any path of symplectic forms on $M$ with
$\omega_0=\omega_c$, such that $D_1$ and $D_2$ remain symplectic
for all $t$. At the fixed point $p$, the $2$-planes
$W_1=T_pD_1$ and $W_2=T_pD_2$ are fixed. Since each $D_i$ is symplectic,
the restriction $\omega_t|_{W_i}$ is nondegenerate for every $t$; hence the
orientation induced by $\omega_t$ on each $W_i$ is constant in $t$.
Choose continuously varying bases $e_1(t),f_1(t)\in W_1$ and
$e_2(t),f_2(t)\in W_2$ with
\[
\omega_t(e_i(t),f_i(t))=1.
\]
Then the orientation of the ordered basis
$(e_1(t),f_1(t),e_2(t),f_2(t))$ of $T_pM$ is independent of $t$.
Since $\omega_t^2$ is a continuous nonvanishing volume form on $T_pM$,
the sign of the Pfaffian of $\omega_t$ in this basis is constant in $t$.
At $t=0$, this Pfaffian is $1-c^2<0$. Thus it is negative for every $t$.

By the lemma, for every $t$ there is no $\omega_t$-tame almost complex
structure on $T_pM$ preserving both $W_1$ and $W_2$.

Finally, suppose an $\omega_t$-tame $J$ with the required local coordinates
existed near $p$. Since two components of $D$ meet at $p$, we would have $r=2$
and local coordinates $(z_1,z_2)$ such that $z_1,z_2$ are $J$-holomorphic and
\[
D_1=\{z_1=0\},\qquad D_2=\{z_2=0\}.
\]
For $X\in W_i=T_pD_i$, we have $dz_i(X)=0$. The condition
\[
dz_i(JX)=i\,dz_i(X)
\]
then gives $dz_i(JX)=0$, so $JX\in W_i$. Thus $J$ would preserve $W_1$ and
$W_2$, contradicting the linear obstruction.

\section{Conclusion}

The answer is \textbf{no}. The torus example above is a closed symplectic manifold
with a symplectic simple normal crossings divisor such that, even after any allowed
deformation of $\omega$ through cohomologous symplectic forms preserving the divisor,
no $\omega$-tame almost complex structure can make the two branches simultaneously
$J$-holomorphic near their intersection.

\begin{remark}
If one imposes the stronger condition that the branches of the divisor are
$\omega$-orthogonal, then the answer is yes: the symplectic Weinstein normal form
gives local Darboux coordinates in which the divisor is a union of coordinate
hyperplanes, and the standard complex structure in those coordinates extends to a
global tame almost complex structure.
\end{remark}

\end{document}
```